\documentclass[12pt,a4paper]{exam}
\usepackage[utf8]{inputenc}
\usepackage[T1]{fontenc}
\usepackage{amsmath}
\usepackage{amsfonts}
%\usepackage{amssymb}
\usepackage{graphicx}
\usepackage{geometry}
\usepackage{enumitem}

\geometry{a4paper, margin=2cm}

\usepackage{cprotect}

\usepackage{xcolor}
\definecolor{maroon}{cmyk}{0, 0.87, 0.68, 0.32}
\definecolor{halfgray}{gray}{0.55}
\definecolor{ipython-frame}{RGB}{207, 207, 207}
\definecolor{ipython-bg}{RGB}{247, 247, 247}
\definecolor{ipython-red}{RGB}{186, 33, 33}
\definecolor{ipython-green}{RGB}{0, 128, 0}
\definecolor{ipython-cyan}{RGB}{64, 128, 128}
\definecolor{ipython-purple}{RGB}{170, 34, 255}

\usepackage{listings}
\lstdefinelanguage{iPython}{
	morekeywords={access,and,del,except,exec,in,is,lambda,not,or,raise},
	morekeywords=[2]{for,print,abs,all,any,basestring,bin,bool,bytearray,callable,chr,classmethod,cmp,compile,complex,delattr,dict,dir,divmod,enumerate,eval,execfile,file,filter,float,format,frozenset,getattr,globals,hasattr,hash,help,hex,id,input,int,isinstance,issubclass,iter,len,list,locals,long,map,max,memoryview,min,next,object,oct,open,ord,pow,property,range,reduce,reload,repr,reversed,round,set,setattr,slice,sorted,staticmethod,str,sum,super,tuple,type,unichr,unicode,vars,xrange,zip,apply,buffer,coerce,intern,elif,else,if,continue,break,while,class,def,return,try,except,import,finally,try,except,from,global,pass, True, False},
	sensitive=true,
	morecomment=[l]\#,%
	morestring=[b]',%
	morestring=[b]",%
	moredelim=**[is][\color{black}]{@@}{@@},
	identifierstyle=\color{black}\footnotesize\ttfamily,
	commentstyle=\color{ipython-cyan}\footnotesize\itshape\ttfamily,
	stringstyle=\color{ipython-red}\footnotesize\ttfamily,
	keepspaces=true,
	showspaces=false,
	showstringspaces=false,
	rulecolor=\color{ipython-frame},
	frame=single,
	frameround={t}{t}{t}{t},
	backgroundcolor=\color{ipython-bg},
	basicstyle=\footnotesize\ttfamily,
	keywordstyle=[2]\color{ipython-green}\bfseries\footnotesize\ttfamily, 
	keywordstyle=\color{ipython-purple}\bfseries\footnotesize\ttfamily
}

\lstdefinelanguage{iOutput} {
	sensitive=true,
	identifierstyle=\color{black}\small\ttfamily,
	stringstyle=\color{ipython-red}\small\ttfamily,
	keepspaces=true,
	showspaces=false,
	showstringspaces=false,
	rulecolor=\color{ipython-frame},
	basicstyle=\small\ttfamily,
}

\lstnewenvironment{ipython}[1][]{\lstset{language=iPython,mathescape=true,escapeinside={*@}{@*}}%
}{%
}

\lstnewenvironment{ioutput}[1][]{\lstset{language=iOutput,mathescape=true,escapeinside={*@}{@*}}%
}{%
}


\title{Financial Market Course 25/26\\ Exam}
\author{Prof. Simone Freschi, Prof. Matteo Sani}
\date{$27^{\mathrm{st}}$ January 2026}

%\printanswers
\noprintanswers
\begin{document}
\maketitle

\begin{center}
\fbox{\fbox{\parbox{5.5in}{
Answer the questions in the provided spaces. If you run out of room for an answer, continue on the page back.}}}
\end{center}

\begin{center}
\vspace{5mm}
\makebox[0.75\textwidth]{Student's name:\enspace\hrulefill}
\end{center}

\section*{Questions}
\vspace{.5cm}
\begin{questions}

\question 
Risk-Performance Evaluation Measures.
\begin{parts}
	\part Describe the Capture Ratio Measures. Imagine that Asset manager \textbf{A} has as an Upside Capture Ratio of 140\% and Downside Capture Ratio of 110\%; Asset manager \textbf{B} has an Upside Capture Ratio of 90\% and Downside Capture Ratio of 70\%; who is the best according to this measure ?
	\part Describe the Draw-Down measure and the Max Draw-Down measure. Imagine that Asset manager \textbf{A} has an expected excess return of 10\% and Max Draw-Down equal to 20\%; Asset manager \textbf{B} has an expected excess return of 5\% and Max Draw-Down equal to 7\%; who would you prefer and why ? Will your choice be different if you had a target return of more or equal to 10\% ?
\end{parts}

\makeemptybox{5cm}
\fillwithdottedlines{5cm}
\begin{solution}
\end{solution}

\question 
Describe what are the main characteristics of a Repo Contract.
\begin{parts}
	\part In chart 1 is described the delivery Basket of the BTP Future March 2025 (IKH5) which is currently trading at 118.63. The bond BTPS 4.35 Nov33 is the cheapest to deliver with a price of 105.50 and a conversion factor of 0.8953. Describe if there is an arbitrage opportunity and what kind of strategy would you implement. (\textit{Hint:} check the basis and use the cash and carry strategy).
\end{parts}

\makeemptybox{5cm}
\fillwithdottedlines{5cm}
\begin{solution}
\end{solution}

%%%%%%%%%%%%%%%%%%%%%%%%%%%%%%%%%%%%%
\question 
Describe the equation for the return attribution model by Brinson, Hood and Beebower (BHB).
\begin{parts}
	\part Explain the difference between Allocation Effect and Security Selection Effect.
	\part Imagine that your portfolio is underweight Tech (10\% vs 20\% Benchmark) and the Tech sector is up 15\%. The Tech stocks in your portfolio are up only 25\%. Determine the allocation effect and the selection effect and the overall portfolio result wrt to sector.
\end{parts}

\makeemptybox{10cm}

\begin{solution}
	\begin{parts}
		\part Calculate the mark-to-market for each scenario:
		\begin{itemize}
			\item A: $1,000\cdot(110-80)=+\$30,000$;
			\item B: $1,000\cdot (80-80)=\$0$;
			\item C: $1,000\cdot (50-80)=-\$30,000$.
		\end{itemize}
		then identify the exposure ($E=\max(MtM,0)$):
		\begin{itemize}
			\item A: \$30,000;
			\item B: \$0;
			\item C: \$0.
		\end{itemize}
		Finally calculate $EE$:
		\begin{equation*}
			EE=\frac{1}{3}(30,000)+\frac{1}{3}(0)+\frac{1}{3}(0)=\$10,000
		\end{equation*}
		
		To calculate the CVA use the formula: $CVA\approx LGD\cdot PD\cdot EE$:
		\begin{align*}
			LGD=1-R=1-0.40=0.60 \\
			CVA=0.60\cdot 0.05\cdot 10,000 \\
			CVA=0.03\cdot 10,000=\$300
		\end{align*}
		\part The adjustment for the bank's own default risk is called DVA (Debit Valuation Adjustment). The total adjustment is BCVA = CVA-DVA.
		\part For an oil producer, Scenario C (MtM=-\$30,000) is when they are likely to default, but the bank's exposure is zero, so no WWR. For it to occur the counterparty should be an oil consumer (like an airline), hence Scenario A would create WWR.
	\end{parts}
\end{solution}

%%%%%%%%%%%%%%%%%%%%%%%%%%%%%%%%%%%%%
\question 
An analyst wants to simulate the short term interest rate using the Vasicek model:
\begin{equation*}
	dr_t = \kappa(\theta - rt)dt + \sigma dW_t
\end{equation*}

with the following parameters:
\begin{itemize}
	\item asymptotic value ($\theta$): 5\%;
	\item mean reversion speed ($\kappa$): 2;
	\item volatility ($\sigma$): 2\%;
	\item initial interest rate ($r_0$) 7\%;
	\item time-horizon ($T$): 0.1 year (with 2 steps: $\Delta t= 0.05$).
\end{itemize}

Solve:
\begin{parts}
	\part using the Euler scheme compute $r_1$ and $r_2$ at the two time steps. The stochastic component can be computed using the following samples from the normal distribution $Z\approx \mathcal{N}(0,1)$:
	\begin{itemize}
		\item Step 1 ($t=0.05$): $z_1=1.0$;
		\item Step 2 ($t=0.10$): $z_2=-1.5$.
	\end{itemize}
	
	\part How would the previous simulation change if the mean reversion speed was very high (e.g. $\kappa = 50$) and $\Delta t=0.05$ ?
	\part Explain why the Euler simulation of a Vasicek model can result in negative interest rates, even when $\theta$ and $r_0$ are positive.
	%\part Suppose to simulate an asset price using a Geometric Brownian Motion (GBM): $dS_t= \mu S_t dt+\sigma S_t dW_t$.
	%Explain why there is necessarily a bias in the simulation using the Euler scheme with respect to the exact solution especially for large $\Delta t$ ?
	%\textbf{Hint}: draw a small piece of the two solutions: exact GBM and its Euler approximation.
\end{parts}
\makeemptybox{20cm}

\begin{solution}
	\begin{parts}
		\part The Euler formula for the Vasicek model is: $r_{t+1}=r_t +\kappa (\theta - r_t)\Delta t+\sigma \Delta t z_{t +1}$
		\begin{itemize}
			\item Solution step 1 ($r_1$):
			\begin{gather*}
				\Delta r= 2(0.05-0.07)(0.05)+(0.02)0.05(1.0) =0.00247 \\
				r_1=0.07+0.00247=7.247\%
			\end{gather*}
			\item Solution step 2 ($r_2$):
			\begin{gather*}
				\Delta r = 2(0.05-0.07247)(0.05)+(0.02)0.05(-1.5) =-0.00895 \\
				r_2=0.07247-0.00895=6.352\%
			\end{gather*}
		\end{itemize}
		\part If $\kappa\Delta t >1$, the "correction" to the short rate is bigger than the distance to the asymptotic value $\theta$, forcing the solution to oscillate around the mean. %If $\kappa \Delta t >2$, Euler scheme becomes unstable and the short rate diverges. 
		\part In the Euler scheme $r_{i+1} = r_i + \kappa(\theta - r_i)\Delta t + \sigma\sqrt{\Delta t} Z$, the term $Z$ is a random sample from a Normal Distribution $\mathcal{N}(0,1)$. Since the Normal distribution has "infinite tails," there is always a non-zero probability of drawing a value of $Z$ so negative that it outweighs the starting rate and the drift.
		%\part The Euler scheme is essentially a linear approximation of the precess (i.e.  in this case an exponential process). Since the exponential is a convex function, the linear approximation will tend to systematically underestimate the expected price of the asset ($\mathbb{E}[S_t]$) in the long run. To compensate for this effect is usually discretized the process $\ln(S_t)$.
	\end{parts}
\end{solution}

%%%%%%%%%%%%%%%%%%%%%%%%%%%%%%%%%%%%%

\question 
Evaluate a Payer Interest Rate Swap (IRS) with a notional $N = € 10_000_000$, maturity 1 year, and annual payments. The IRS fixed rate is $K = 3\%$.
At $t=0$, the 1y-EURIBOR rate and the risk free rate are both $4\%$. Briefly explain the result.
\makeemptybox{10cm}

\begin{solution}
A payer IRS value can be seen as the difference between a variable rate bond and a fixed rate bond:
$NPV = V_{float} - V_{fix}$

Since the EURIBOR rate $L$ is equal to the discounting rate $r$ the floating leg equals par:
$V_{float} = \frac{N(1 + L)}{1 + r} = \frac{10000000 \times 1.04}{1.04} = 10000000$

The fixed leg instead:
$V_{fix} = \frac{N(1 + K)}{1 + r} = \frac{10000000 \times 1.03}{1.04} = 9903846$

Hence the NPV:
$NPV = 10000000 - 9903846 = 96154$

Since interest rates increased the party paying fixed has a gain, she pays 3\% and receives 4\%.
\end{solution}

\end{questions}
\end{document}
